\begin{claim}{C.87}{Proved}{Global contractions of the radial and spiral relative centralizers}
The entire topological groups $\mathscr P$ of
C.84, $\mathscr C_{\mathrm{rad}}$ of C.85,
and $\mathscr C_{\mathrm{sp}}$ of C.86
are contractible by continuous contractions
fixing the identity throughout.
The proof constructs one contraction on
the whole group from the parametric C.85
cutoff, a continuous support height, an
exact compact-annulus lift comparison,
and the compact boundary-relative
identity-component contraction.
Every original disk boundary point and
every specified integer operator label
is retained, including in the inverse
topology. The compact-annulus comparison
identifies its identity component exactly
with the kernel of the displayed integer
boundary-shift homomorphism.
\end{claim}

\begin{proof}
\emph{The exact compact-annulus comparison.}
Put $A=(\mathbb R/\mathbb Z)\times[-1,1]$
and let $\mathscr H$ be its
orientation-preserving homeomorphism
group fixing both boundary circles
pointwise. Each $b\in\mathscr H$
induces the identity on the fundamental
group, since the lower boundary is
fixed. It has a unique lift
$\widetilde b:\mathbb R\times[-1,1]
\to\mathbb R\times[-1,1]$ fixing
$(0,-1)$. Its restriction to the
whole lower boundary is the identity,
and uniqueness of lifts gives
$\widetilde bT_1=T_1\widetilde b$.
On the upper boundary it has the form
\[
 \widetilde b(x,-1)=(x,-1),\qquad
 \widetilde b(x,1)=(x+\tau(b),1),\qquad
                         \tau(b)\in\mathbb Z.                   \tag{C.87.1}
\]
The integer is independent of $x$
by connectedness. The canonical lift
of a composite is the composite of
the canonical lifts, so $\tau$ is
a homomorphism. The maps
\[
     b_k([x],y)=\bigl([x+k(y+1)/2],y\bigr)
                         \quad(k\in\mathbb Z)                  \tag{C.87.2}
\]
fix both boundary circles and have
$\tau(b_k)=k$; they give a literal
homomorphic section.

The canonical lift varies continuously
in the compact-open map-and-inverse
topology, uniformly on one period
rectangle. To verify this directly,
let $b_i\to b$. For large $i$, the
two circle coordinates of $b_i(z)$
and $b(z)$ differ by a unique small
real value $\delta_i(z)\in(-1/4,1/4)$,
continuous in $z\in A$ and tending
uniformly to zero. The difference
of their lifted first coordinates
minus $\delta_i$ is a continuous
integer on the connected covering
strip. It is zero on the lower
boundary and hence zero everywhere.
Their second coordinates converge
uniformly as well. Apply the same
argument to inverses. In particular
$\tau(b_i)=\tau(b)$ eventually, since
$\delta_i=0$ on the upper boundary.
Thus $\tau$ is continuous with
discrete target.

We now prove its exact kernel.
Suppose $\tau(b)=0$ and let
$\gamma(y)=([0],y)$ be the spanning
arc. Its image has a lift
$\widetilde b(0,y)$ with endpoints
$(0,-1)$ and $(0,1)$.
The straight homotopy in the convex
strip from this lifted arc to
$(0,y)$ fixes these two endpoints;
its interior stays in the interior
of the strip. Its projection proves
that $b\gamma$ and $\gamma$ are
homotopic with their endpoints
fixed. Epstein's actual relative
arc theorem
\cite[Theorem~3.1, pp.~89--90; Appendix~A1--A2 and their proofs, pp.~105--106]{C_Epstein1966}
gives an ambient isotopy fixed
pointwise on both boundary circles
carrying $b\gamma$ to $\gamma$.
The Appendix supplies the stated
topological arc case: A2's proof
first makes the arc piecewise linear
near its endpoints and then applies
the construction in the proof of A1.

Any remaining parametrization of
$\gamma$ is an increasing interval
homeomorphism $\rho$ fixing its
endpoints. In a rectangular
neighborhood with coordinates
$(u,y)$ and $\gamma=\{u=0\}$,
choose $0\leq\zeta\leq1$ with
$\zeta(0)=1$ and $\zeta=0$ on
the two side edges. The formula
\[
       (u,y)\longmapsto
       \bigl(u,(1-s\zeta(u))y+s\zeta(u)\rho^{-1}(y)\bigr)
                              \quad(0\leq s\leq1)               \tag{C.87.3}
\]
is an isotopy of that rectangle,
fixed on its entire boundary,
and has a continuous inverse by
monotone bracketing. Extend it by
the identity. After this correction
the composed annulus map fixes
$\gamma$ pointwise. Cutting along
$\gamma$ produces a disk map
fixing its boundary pointwise.
The disk Alexander isotopy descends
to an isotopy of the annulus
relative to its boundary and
$\gamma$. Thus $b$ belongs to the
identity path component of
$\mathscr H$. Conversely the
continuous integer $\tau$ vanishes
on the identity component. We have
proved the exact equality
\[
                         \ker\tau=\mathscr H_0.                 \tag{C.87.4}
\]
The arc isotopies in this argument
establish membership in that
component; no parameter choice
of those isotopies is used below.

Let
\[
 \mathscr E=\{F\in\mathscr P:
                    F(x,y)=(x,y)\text{ for }|y|\geq1\}.
                                                               \tag{C.87.5}
\]
Restriction and descent identify
$\mathscr E$ with $\mathscr H_0$.
The actual inverse map
$\Lambda:\mathscr H_0\to\mathscr E$
takes the lower-fixed lift and
extends it by the identity outside
the strip. Equation (C.87.4) fixes
the upper line as well, so this
pasting is an actual plane
homeomorphism. Its inverse is
$\Lambda(b^{-1})$ and its
translation labels are unchanged.
Its displacement is bounded and
periodic, so it lies in
$\mathscr P$.
These are inverse topological
group maps. Restriction and
descent are continuous on the
compact period rectangle by
C.84. The canonical lift
continuity just proved, together
with its identity extension,
gives continuity of $\Lambda$
in both weighted terms of
(C.84.2).

The compact boundary-relative
identity component is contractible
by
\cite[Lemma~2.1(ii), pp.~2--3]{C_Yagasaki2000}.
Here the whole nonempty boundary
of the compact annulus is fixed
pointwise, as required there.
Choose one continuous contraction
$C_s$ of $\mathscr H_0$ with
$C_0(b)=b$ and $C_1$ constant.
Replacing it by
$C_s(b)C_s(1)^{-1}$ fixes the
identity throughout and makes
the last map the constant
identity. Transport it through
$\Lambda$ and denote the resulting
single continuous contraction
on the whole of $\mathscr E$ by
\[
 \widehat\Gamma:[0,1]\times\mathscr E\to\mathscr E,
 \quad
 \widehat\Gamma_0(G)=G,\quad
 \widehat\Gamma_1(G)=1,\quad
 \widehat\Gamma_s(1)=1.
                                                               \tag{C.87.6}
\]

\emph{The actual vertical normalizers.}
For $w>0$ put $D_w(x,y)=(x,wy)$
and
\[
 \begin{gathered}
 \mathcal N_w(F)=D_wFD_w^{-1},\\
 \mathcal N_w(F)(x,y)
       =\bigl(X_F(x,y/w),\ wY_F(x,y/w)\bigr).
 \end{gathered}                                                   \tag{C.87.7}
\]
The exact inverse parameter is
$w^{-1}$ and $D_w$ commutes with
the original $T_1$.
Sublinearity in (C.84.1) is
preserved by the displayed
coordinate formula. For fixed
$w$ the weighted metric satisfies
\[
 \delta_*(\mathcal N_w(F),\mathcal N_w(H))
                  \leq\max\{w,w^{-1}\}\,\delta_*(F,H).
                                                               \tag{C.87.8}
\]
Indeed the Euclidean operator
norm of $D_w$ is $\max\{1,w\}$,
and the change of source height
has ratio
$(1+|y|/w)/(1+|y|)\leq\max\{1,w^{-1}\}$.
Apply the same estimate to the
exact inverse conjugations of
$F^{-1}$ and $H^{-1}$.

For joint continuity in $w$
retain the original closed disk.
For $0\leq\rho\leq1$ and $|u|=1$
the actual normalizer is
\[
 \overline D_w(\rho u)
    =\frac{\rho D_wu}{1-\rho+\rho|D_wu|},\qquad
 \overline D_w(0)=0,\qquad
 \overline D_w^{-1}=\overline D_{1/w}.
                                                               \tag{C.87.9}
\]
For $\rho<1$ it is exactly
$jD_wj^{-1}$, and on the circle
it is $D_wu/|D_wu|$.
On any compact positive interval
of $w$, the denominator is
bounded away from zero. The
formula is jointly continuous,
also at $\rho=0$, and so is the
displayed inverse. It gives
joint continuity into the
closed-disk homeomorphism group.
It commutes with the specified
operator $jT_1j^{-1}$. Conjugation
by it preserves pointwise
boundary fixation. Transporting
that conjugation through the
exact C.84 isomorphism proves
joint continuity of (C.87.7)
as a map $(0,\infty)\times
\mathscr P\to\mathscr P$,
with both original inverse
topology maps included.

\emph{A single contraction on the entire group.}
Let $L_+(F)=L(F,1)$ be the
positive cutoff height of
(C.85.3), and let
$L_-(F)=L(JFJ,1)$ use the
reflection $J$ of C.85.
Both are continuous functions
of $F$. Put
\[
 \begin{gathered}
 W(F)=3+2\max\{L_+(F),L_-(F)\},\\
 G(F)=\mathcal N_{1/W(F)}(\mathcal C_1(F)).
 \end{gathered}                                                   \tag{C.87.10}
\]
The cutoff is the identity beyond
heights $1+2L_+(F)$ and
$-1-2L_-(F)$. Since each of their
absolute values is at least two
less than $W(F)$, the map
$G(F)$ is the identity near
both boundary lines of the
fixed strip and outside it.
Thus $G(F)\in\mathscr E$.
It depends continuously on
$F$ by C.85 and (C.87.7).
The exact original lift fixes
both boundary lines, so the
annulus component used in
(C.87.6) is indeed the
zero-$\tau$ component just
proved, for every $F$.

Define on the whole group
\[
 \mathcal H_s(F)=
 \begin{cases}
  F,&s=0,\\
  \mathcal C_{1/(2s)}(F),&0<s\leq\tfrac12,\\
  \mathcal N_{W(F)}
    \bigl(\widehat\Gamma_{2s-1}(G(F))\bigr),
       &\tfrac12\leq s\leq1.
 \end{cases}                                                     \tag{C.87.11}
\]
C.85 gives joint continuity
of the first part, including
at zero. At $s=1/2$ the second
part equals
$\mathcal N_{W(F)}G(F)=\mathcal C_1(F)$,
so the two parts agree exactly.
Every component map in the
last expression is continuous
on its full stated domain.
The function $W$ is locally
bounded and positive; in
particular the last part near
any fixed $F$ has support in
one common horizontal band,
and its inverse has that
same support. The normalizer
comparison already proves
the required original
map-and-inverse continuity.
Pasting gives a single
continuous map
$[0,1]\times\mathscr P\to\mathscr P$.
Its faces are
\[
 \mathcal H_0(F)=F,\qquad
 \mathcal H_1(F)=1,\qquad
 \mathcal H_s(1)=1.
                                                               \tag{C.87.12}
\]
For the last equality use
$\mathcal C_R(1)=1$ and
$\widehat\Gamma_s(1)=1$;
no normalization of $W(1)$
is required.

Conjugate (C.87.11) through
the C.84 topological group
isomorphism to obtain the
contraction on
$\mathscr C_{\mathrm{rad}}$.
Conjugate that contraction
through the two actual
topology maps $\Theta$ and
$\Theta^{-1}$ of C.86 to
obtain the contraction on
$\mathscr C_{\mathrm{sp}}$.
These formulas act on the
whole respective groups,
fix the identity, and
preserve the original
boundary and integer
operator labels throughout.
The general discrete-operator
case and general C.12 remain
open.
\end{proof}
